\documentclass[10pt,a4paper]{amsart}
\usepackage[margin=19mm]{geometry}
\usepackage{amsmath,amssymb,mathtools}
\usepackage[scr=rsfs,cal=euler]{mathalfa}
\usepackage{xfrac}
\usepackage[unicode,hidelinks,bookmarks=false]{hyperref}
\newcommand{\ie}{i.e.\ }
\newcommand{\cf}{cf.\ }
\newcommand{\resp}{respectively\ }
\newcommand{\Subsection}{\subsection}
\providecommand{\nobreakdash}{\nobreak\hbox{-}}
\setlength{\emergencystretch}{2em}
\multlinegap0pt
\usepackage{bm}
\usepackage{bbm}
\usepackage{dsfont}
\usepackage{xspace}
\usepackage{cancel}
\usepackage{mathtools}
\usepackage{amsthm}
\makeatletter
\@ifundefined{lemma}{\newtheorem{lemma}{Lemma}}{}
\makeatother
\newcommand{\m}{\mathcal}
\title[Families without $s$-matchings: the other end]{Families without $s$-matchings: the other end}
\author{Andrey Kupavskii, Georgy Sokolov}
\date{Source: arXiv:2605.00996v1 (CC BY 4.0)}
\begin{document}
\maketitle

\noindent\emph{Source and licence.} Families without $s$-matchings: the other end. Version arXiv:2605.00996v1 (CC BY 4.0), distributed under the Creative Commons Attribution 4.0 International licence, as recorded in the arXiv licence metadata for this exact version and shown on the abstract page https://arxiv.org/abs/2605.00996v1. Full rights notice: licenses/arxiv-2605.00996-CC-BY-4.0.txt.\par\smallskip

\noindent\textit{Editorial scope.} Counted unit: the Lemma whose source label is \texttt{l.stability\_d\_eq\_mc-c} in the article (source0.tex lines 387--392, source label \texttt{l.stability\_d\_eq\_mc-c}), delivered together with the article's own complete original proof of it (source0.tex lines 394--479, one \texttt{proof} environment of 86 lines). No mathematics is written, repaired or re-derived: statement and proof are the article's own text line for line. Required context retained uncounted: l.deleteon\_with\_shifting (lines 258--260); l.d\_geq\_mc-c+1 (lines 340--345). Every definition, display and label of those lines is retained unchanged. Omitted from this package and not redistributed: 1--257, 261--339, 346--386, 393--393, 480--622. Nothing inside a retained excerpt is omitted or altered: every retained body is delivered exactly as the article's lines are written, so no omission receipt of any kind accompanies this package. Citations: the retained excerpts carry no citation command at all, so no bibliography entry of the article is transcribed and no reference is redistributed. Reference adaptation: none was needed and none was made; every reference inside the delivered unit resolves inside it, so no unresolved reference or citation is produced. Printed slips kept literal: no printed word, symbol, hypothesis or number of the counted statement or of its proof was altered. Layout and class: the article's own class is \texttt{article} with packages including inputenc, fontenc, amscd, amssymb, amsfonts, amsmath, array, graphicx, longtable, wrapfig, amsthm, xcolor; the wrapper is the lane's pinned \texttt{amsart} editorial wrapper at 10pt on A4, which restates the theorem-like environments the retained text needs and adds the article's own macro definitions that the retained text needs (\texttt{\textbackslash m}), with extra wrapper packages (bm, bbm, dsfont, xspace, cancel, mathtools). The delivered numbering is the unit's own. Assets: no retained span contains a figure environment, a graphics-inclusion command, a PDF-page inclusion, a table or a TikZ picture; no image file is copied into this package, the package declares no asset dependency and no image file is redistributed with it. Licence: version arXiv:2605.00996v1 of this article is distributed under the Creative Commons Attribution 4.0 International licence, as recorded in the arXiv licence metadata for this exact version and shown on the abstract page https://arxiv.org/abs/2605.00996v1. A full rights notice is delivered at licenses/arxiv-2605.00996-CC-BY-4.0.txt. Characters: every retained excerpt is pure ASCII, so the delivered engine, XeLaTeX with its default Latin Modern text font, reports no missing character.
\begin{lemma} \label{l.deleteon_with_shifting}
    If $\mathcal{F}$ is shifted, $\ell \geq 1$, and $d(\mathcal{F}) > d$, then for any integer $k$, $0 \leq k \leq \frac{d}{m}$, and any $X \subset [d + m\ell]$ with $|X| = d - km$, the family $\mathcal{F} \cap {[d+m\ell] \setminus X \choose m}$ contains an $(\ell+k)$-matching.
\end{lemma}

\begin{lemma} \label{l.d_geq_mc-c+1}
    If $\mathcal{F} \subset 2^{[n]}$ is a shifted family with $d(\mathcal{F}) > (m-1)c$, then
    \[
    y_{\mathcal{F}}(m+1) \geq (2c-1){n \choose m - 1} + O(s^{m - 2}).
    \]
\end{lemma}

\begin{lemma} \label{l.stability_d_eq_mc-c}
    If $\mathcal{F} \subset 2^{[n]}$ is a shifted family with $d(\mathcal{F}) = (m-1)c$ and $\mathcal{F} \nsubseteq \mathcal{Q}(m,s,\ell)$, then
    \[
    y_{\mathcal{F}}(m+1) \geq {n \choose m - 1} + O(s^{m-2}).
    \]
\end{lemma}

\begin{proof}
    The proof is very similar to that of Lemma~\ref{l.d_geq_mc-c+1}, although the construction of a random $s$-matching is slightly more complicated. Let $F$ be an arbitrary set in $\mathcal{F} \setminus \mathcal{Q}(m, s, \ell)$. Denote
    \[
    x = |F \cap [ms-c-1]| \qquad \text{and} \qquad y = |F \setminus [ms-c-1]|.
    \]
    By the definition of $\mathcal{Q}(m, s, \ell)$, we have \begin{equation}\label{eqxy}2x + y \leq 2m - 1.\end{equation}

    Assume that $x \leq m-c-1$. Since $x \geq 0$, we have $m \geq c + 1$, and therefore
    \[
    d(\mathcal{F}) = (m-1)c > mc-c - 1 \geq m(c-1).
    \]
    Thus, we may apply Lemma~\ref{l.deleteon_with_shifting} with $d = mc-c-1$ and $k=c-1$, and conclude that for any
    $
    X \in {[ms-c-1] \choose m-c-1}
    $
    the family $\mathcal{F}\cap {[ms-c-1]\setminus X \choose m}$ contains an $(s-1)$-matching. Since
    \[
    x = |F \cap [ms-c-1]| \leq m - c - 1,
    \]
    it follows that $\mathcal{F}\cap {[ms-c-1]\setminus F \choose m}$ contains an $(s-1)$-matching. Adding $F$ to this matching, we obtain an $s$-matching in $\mathcal{F}$, a contradiction with $\nu(\m F)<s$.

    Thus, we may assume that $x \geq m - c$. Construct a random $s$-matching $\pi$ in
    \[
    \Bigl(\mathcal{F} \cap {[ms-c-1] \choose m}\Bigr) \cup \Bigl({[ms-c-1] \choose m-1} \times {[ms-c, n] \choose 2}\Bigr) \cup \{F\}
    \]
    by the following procedure.
    \begin{itemize}
        \item Add $F$ to $\pi$.
        \item Choose
        $
        X \in {[ms-c-1] \setminus F \choose (m-1)(x+c+1-m)}
        $
        uniformly at random. Note that $x+c+1-m > 0$, since $x \geq m - c$.
        \item By Lemma~\ref{l.deleteon_with_shifting}, with $d = mc-c-1$ and $k = m-x-2$, for any
        $
        Y \in {[ms-c-1] \choose d-km}
        $
        the family $\mathcal{F} \cap {[ms-c-1] \setminus Y \choose m}$ contains an $(\ell+m-x-2)$-matching. We have $d-km=(m-1)(x+c+1-m) + x$. We apply the lemma to
        $
        Y = X \cup (F \cap [ms-c-1])
        $
        and obtain an $(\ell+m-x-2)$-matching in
        \[
        \mathcal{F} \cap {[ms-c-1] \setminus (X \cup F) \choose m}.
        \]
        Add this matching to $\pi$.
        \item Choose a random $(x+c+1-m)$-matching in
        \[
        {X \choose m-1} \times {[ms-c, n] \setminus F \choose 2}
        \]
        and add it to $\pi$. Note that such matchings exist, since $|X|=(m-1)(x+c+1-m)$ and
        \[
        |[ms-c, n] \setminus F| = 2c+1-y \overset{\eqref{eqxy}}{\geq} 2(x+c+1-m).
        \]
    \end{itemize}

    Note that $\pi$ contains
    \[
    1+(\ell+m-x-2)+(x+c+1-m) = s
    \]
    sets, and all sets except those added in the last step belong to $\mathcal{F}$. Let $\xi$ be the number of sets from $\overline{\mathcal{F}}$ in $\pi$. On the one hand, $\xi \geq 1$, since $\pi$ is an $s$-matching and $\mathcal{F}$ does not contain an $s$-matching. Therefore, $\mathbb{E}\xi \geq 1$. On the other hand, for any set
    \[
    H \in \m G:={[ms-c-1] \setminus F \choose m-1} \times {[ms-c, n] \setminus F \choose 2}
    \]
    we have
    \[
    \mathbb{P}[H \in \pi] = (x+c+1-m){ms-c-1-x \choose m-1}^{-1}{2c+1-y \choose 2}^{-1}.
    \]
    Therefore,
    \[
    1\le \mathbb{E}\xi = |\overline{\mathcal{F}} \cap \m G|
    \cdot (x+c+1-m){ms-c-1-x \choose m-1}^{-1}{2c+1-y \choose 2}^{-1}.
    \]
    Combining these inequalities, we get
    \[
    \begin{split}
        y_{\mathcal{F}}(m+1)
        &\geq |\overline{\mathcal{F}} \cap \m G| \\
        &\geq (x+c+1-m)^{-1}{ms-c-1-x \choose m-1}{2c+1-y \choose 2} \\
        &\geq (x+c+1-m)^{-1}{ms-c-1-x \choose m-1}{2(x+c+1-m) \choose 2} \\
        &= (2x+2c+1-2m){n-2c-1-x \choose m-1} \\
        &\geq {n-2c-1-x \choose m-1} \\
        &= {n \choose m - 1} + O(s^{m - 2}).
    \end{split}
    \]
\end{proof}

\end{document}
